% ============================================================================
%  Section I — Introduction  (IOP MST; passive voice / third person)
% ============================================================================

\section{Introduction}
\label{sec:intro}

Induction machines account for a dominant share of the electrical
energy consumed by industry and form the mechanical backbone of countless
processes; their unexpected failure causes costly downtime and, in critical
applications, safety hazards. Continuous condition monitoring (CM) and early
fault diagnosis are therefore central to modern maintenance
practice~\cite{gangsar2020signal, hamani2025advancements}. Among the available
modalities, vibration analysis is the most established, because the mechanical
and electrical faults of interest---bearing defects, rotor asymmetries, and
supply imbalances---imprint characteristic signatures on the machine's vibration
response~\cite{yang2021fault, zhang2020deep}.

The prevailing paradigm relies on industrial-grade piezoelectric accelerometers
sampled at tens to hundreds of kilohertz, which resolve the high-frequency
bearing defect frequencies exploited by spectral and envelope analysis. This is
reflected in the public benchmarks that anchor the field, acquired at
$8$--$200$~kHz~\cite{lee2024multidomain, sehri2024ottawa, thuan2023hust,
huang2018bearing}. Although accurate, such instrumentation is costly, requires
intrusive installation, and scales poorly to the vast population of small,
low-value machines that individually do not justify a dedicated monitoring channel
yet collectively dominate a plant's machine inventory. Closing this coverage gap
motivates a shift toward inexpensive, wireless, and readily deployable sensing.

Consumer micro-electro-mechanical-systems (MEMS) inertial sensors, epitomised by
those embedded in smartphones, offer such a platform at negligible marginal
cost~\cite{grouios2022accelerometers}. Machine fault detection has been demonstrated
from smartphone- and MEMS-acquired vibration~\cite{naha2017mobile,
ertargin2023deep, ertargin2025automated}, and dedicated low-cost MEMS nodes now
integrate on-board processing with wireless telemetry for machine
diagnostics~\cite{pedotti2020lowcost, jakobsen2024lowcost}. These sensors,
however, operate at sampling rates one to three orders of magnitude below
industrial practice---$100$~Hz in the present case---so that, by the Nyquist
criterion~\cite{nyquist1924certain, shannon1948mathematical}, the high-frequency
signatures underpinning classical bearing diagnosis are not observable. A
fundamental and largely unresolved question follows: what fault information is
genuinely recoverable from such a constrained sensor, and does the apparent
performance transfer to a real deployment?

The second half of that question is easily overlooked. Learning-based fault
diagnosis routinely reports near-perfect accuracy, yet a growing body of evidence
attributes much of this success to \emph{data leakage}---information shared
between training and test partitions that inflates measured accuracy without
improving generalisation~\cite{kapoor2023leakage}. In vibration diagnosis, leakage
arises when overlapping analysis windows from a single continuous recording are
divided across partitions, or when recording-specific characteristics are learned
in place of fault features; enforcing component- or recording-wise separation has
been shown to reduce reported accuracy substantially~\cite{wheat2024impact,
hendriks2022towards}. This scrutiny has so far been directed at industrial-grade,
high-rate data. For low-rate smartphone datasets---in which each operating
condition is captured in one continuous acquisition and heavy window overlap is
common---the risk is arguably greater, yet random-split evaluation remains the
norm, including in the baseline accompanying the dataset examined
here~\cite{ertargin2026smartphone}.

A third consideration is deployment. The Industrial Internet of Things envisions
dense networks of inexpensive nodes that stream, or locally distil, diagnostic
information over constrained wireless links~\cite{yousuf2024iot}, and recent
TinyML implementations have shown that lightweight classifiers can run on
microcontrollers at millisecond latency and sub-millijoule
energy~\cite{gao2025edge, asutkar2023tinyml}. Such work establishes that edge
diagnosis is feasible, but it generally presupposes industrial-bandwidth sensing
and does not characterise the coupled trade-off between diagnostic accuracy and
the sampling, communication, and computational cost that governs an ultra-low-rate
node. Whether, for a smartphone-class sensor, the low sampling rate is a liability
to be tolerated or an asset to be exploited has not been quantified.

This paper approaches these gaps as a measurement problem. Using a public
smartphone-acquired vibration dataset of an induction machine recorded across six
health states, three supply frequencies and two load conditions at
$100$~Hz~\cite{ertargin2026smartphone, ertargin2026dataset}, and rather than
proposing a new classifier, the study treats diagnostic performance as a quantity
to be measured: a protocol is developed for estimating it without leakage, the
mechanisms that inflate it are isolated experimentally, the accuracy--resource
trade-off that governs an IoT node is mapped, and every figure is reported with an
interval whose sampling unit is stated. The scope is deliberately bounded, and the
bound is stated at the outset: all recordings come from one $1.1$~kW machine in
which each fault class is represented by a single physical specimen, so what is
estimated is performance on an unseen recording of that machine, not performance
on an unseen machine.

The principal contributions are fourfold. First, the study \emph{decomposes} the
optimism of conventional evaluation into its mechanisms. Random splitting of
windows is shown to overstate accuracy by $41$ to $47$ percentage points, and a
ladder of partitionings then attributes that gap: window overlap, the mechanism
most often blamed and most often guarded against, accounts for none of it, as does
simple temporal proximity, while recording identity accounts for all $46.7$
points. Enforcing non-overlapping windows, a common precaution, is therefore no
protection at all on data of this kind. Second, it resolves what a sampling-rate
sweep actually measures. Decimation to $25$~Hz improves leakage-free accuracy even
though it removes every shaft rotational frequency in the dataset, and a
controlled comparison shows that band-limiting alone is mildly harmful, that the
improvement follows from the reduced number of samples per window, and that
omitting the anti-alias filter is severely damaging; a rate sweep performed with a
fixed time-domain feature set thus confounds the observable band with the
sample-count sensitivity of the features. Third, it provides a data-derived
characterisation of the smartphone MEMS channel---its observed output increment,
an in-operation signal floor, the band made observable by the sampling rate, and
the relative informativeness of the raw and gravity-compensated triads---while
stating explicitly that these are properties of the delivered signal and not a
metrological calibration of the transducer. Finally, it maps the
accuracy--resource trade-off into deployable terms through a payload-rate
comparison and an edge-model footprint, and reports all comparisons with
recording-clustered intervals and paired permutation tests rather than with
statistics computed over repeated re-partitionings of the same data.

The remainder of this paper is organised as follows.
Section~\ref{sec:related} reviews related work on vibration-based condition
monitoring, low-cost MEMS sensing, IoT and edge diagnosis, and evaluation
methodology. Section~\ref{sec:dataset} describes the dataset and measurement
setup, and Section~\ref{sec:charact} characterises the sensor and its signals.
Section~\ref{sec:method} details the processing chain and the evaluation
protocols. Section~\ref{sec:results} reports the results,
Section~\ref{sec:discussion} discusses their implications and limitations, and
Section~\ref{sec:conclusion} concludes.
